\section{Optimasi Variasional dan Evaluasi Performa}
Proses optimasi dalam penelitian ini dilakukan melalui skema loop hibrida kuantum-klasik, di mana komputer kuantum bertugas menghitung nilai ekspektasi energi dari Hamiltonian sementara komputer klasik memperbarui seluruh parameter sirkuit. Algoritma \textit{Simultaneous Perturbation Stochastic Approximation} (SPSA) dipilih sebagai optimizer klasik karena efisiensinya yang tinggi dalam menangani lanskap energi yang memiliki derau (\textit{noisy landscapes}) pada hardware era NISQ. Selama proses pelatihan yang iteratif, parameter rotasi pada sirkuit kuantum terus diperbarui secara dinamis hingga mencapai titik konvergensi di mana nilai energi Hamiltonian mencapai level minimum yang paling stabil. Penurunan nilai energi secara bertahap ini secara teknis mencerminkan peningkatan kualitas portofolio dalam memitigasi risiko sistemik sekaligus memaksimalkan utilitas strategis berdasarkan data pasar yang telah diolah sebelumnya.

Simulasi \textit{backtesting} dijalankan secara kronologis menggunakan data historis saham LQ45 untuk menguji keandalan serta ketangguhan model dalam menghadapi perubahan rezim pasar yang bersifat dinamis dan tidak menentu. Strategi penyeimbangan portofolio (\textit{rebalancing}) dilakukan secara periodik, di mana sirkuit kuantum akan dilatih ulang menggunakan data dari jendela waktu terbaru guna menangkap pergeseran struktur korelasi antar aset. Performa dari portofolio yang dihasilkan dievaluasi secara ketat menggunakan metrik standar industri yang meliputi \textit{Total Return}, \textit{Sharpe Ratio}, serta \textit{Maximum Drawdown} (MDD) guna mendapatkan gambaran profil risiko-imbal hasil yang utuh. \textit{Sharpe Ratio} dimanfaatkan untuk mengukur efisiensi portofolio dalam menghasilkan imbal hasil tambahan untuk setiap unit risiko yang diambil oleh investor di tengah fluktuasi harga pasar yang stokastik. Perbandingan performa antara model HE-VQE dengan strategi \textit{benchmark} konvensional dilakukan untuk memvalidasi secara empiris keunggulan serta nilai tambah dari penggunaan teknologi komputasi kuantum dalam optimasi finansial.
